%% suggestions.tex
%% Suggested edits/additions to the Methods section of oja_template.tex,
%% based on comparison with:
%%   - FORGE'd in FIRE (Hopkins et al. 2023, arXiv:2309.13115)
%%   - FORGE'd in the Early Universe (Meziani et al. 2026, arXiv:2602.02953)
%%   - FIRE-2 (Hopkins et al. 2018) and FIRE-3 (Hopkins et al. 2022)
%%
%% Each item below gives: (a) what is missing/wrong, (b) suggested text.
%% Items marked [CRITICAL] need author input. Items marked [ADD] can be
%% inserted directly. Items marked [VERIFY] need cross-checking with the
%% actual simulation setup.

%% ============================================================
%% 1.  SECTION INTRO: add explicit reference to Meziani+2026
%% ============================================================
%%
%% CURRENT:
%%   "we utilize FORGE'd in FIRE framework, similar to \citep{FORGE_d_in_FIRE_Hopkin_2023}"
%%
%% ISSUE: Meziani+2026 is the direct methodological predecessor for the
%%   Pop III application of this framework. Both papers should be cited.
%%   Also "A lot of parameters" is informal.
%%
%% SUGGESTED REPLACEMENT:

For this paper we utilise the FORGE'd in FIRE framework of
\citet{FORGE_d_in_FIRE_Hopkin_2023}, extended to the primordial
star-formation regime following \citet{meziani26}.
We employ FIRE-3 \citep{FIRE_3_Updated_Hopkin} physics modules at
low resolution and STARFORGE \citep{STARFORGE_Towa_Grudic_2021}
physics to resolve individual protostellar formation.
Numerical parameters not specified below are identical to those
described in those references.


%% ============================================================
%% 2.  HYDRODYNAMICS: add constrained-gradient MHD citation
%% ============================================================
%%
%% CURRENT cites Hopkins_Raives_2015 for MFM.
%% ISSUE: the constrained-gradient divergence-cleaning method is a
%%   separate paper (Hopkins 2016, MNRAS 462, 576) and is what FORGE
%%   specifically uses for greater accuracy ("using the constrained-gradient
%%   method from \citet{hopkins:cg.mhd.gizmo} for greater accuracy").
%%   This should be added; it also enables accurate MRI capture.
%%
%% ADD to the MHD sentence:

The ideal MHD equations are integrated as described in
\citet{Hopkins_Raives_2015}, using the constrained-gradient method
\citep{Hopkins_2016_CG_MHD} for divergence control, which has been
validated for magnetorotational instability dynamics across a range
of disk geometries \citep{FORGE_d_in_FIRE_Hopkin_2023}.

%% NEW BIB ENTRY NEEDED:
%% @article{Hopkins_2016_CG_MHD,
%%     author = {Hopkins, Philip F.},
%%     title  = "{Constrained-gradient MHD in GIZMO}",
%%     journal = {\mnras},
%%     volume  = {462},
%%     pages   = {576--587},
%%     year    = {2016},
%%     doi     = {10.1093/mnras/stw1578}
%% }


%% ============================================================
%% 3.  NON-IDEAL MHD: add justification for using ideal MHD
%% ============================================================
%%
%% CURRENT mentions non-ideal terms (ambipolar, Hall, Ohmic) but does not
%%   explain why they are sub-dominant in this simulation.
%% ISSUE: Meziani+2026 explicitly validates that for metal-free gas at
%%   these temperatures and densities, NIMHD timescales exceed dynamical
%%   timescales by factors of 10^6-10^10, because of high temperatures,
%%   low densities, and elevated free-electron fractions. This should be
%%   stated to justify using ideal MHD as the dominant regime.
%%
%% ADD after the non-ideal MHD sentence:

Although non-ideal MHD terms (ambipolar diffusion, Hall effect,
Ohmic resistivity) are included via the implementation of
\citet{Hopkins_2017}, for metal-free gas at the temperatures
($T \sim 200$--$10^{4}$~K) and free-electron fractions
($x_e \gg 10^{-4}$) relevant here, the non-ideal timescales
exceed the local dynamical timescale by many orders of magnitude
\citep{meziani26}, so ideal MHD dominates throughout the
simulation.


%% ============================================================
%% 4.  MHD: initial magnetic field value
%% ============================================================
%%
%% CURRENT: no mention of initial field strength.
%% ISSUE: Meziani+2026 states they start from trace cosmological fields
%%   of ~10^-15 G, amplified self-consistently. This should be stated.
%%   [CRITICAL] Confirm your initial B-field value.
%%
%% SUGGESTED:

The initial magnetic field is a trace cosmological seed field of
$|B_0| \sim 10^{-15}$~G \note{[VERIFY: confirm initial B value]},
amplified self-consistently by gravitational compression,
differential rotation, and the turbulent dynamo throughout the
simulation.


%% ============================================================
%% 5.  RADIATION: add reduced speed of light, clarify bands
%% ============================================================
%%
%% CURRENT: lists 7 radiation bands (photo-ionizing, LW, photo-electric,
%%   NUV, optical, NIR, FIR) but FORGE uses 5 bands where optical+NIR
%%   are combined and FIR is adaptive-wavelength.
%%   Missing: reduced speed of light value.
%%   Missing: reference for M1 implementation in GIZMO (Hopkins radiation
%%   methods paper, Lupi et al.).
%%
%% SUGGESTED REPLACEMENT for the Radiation paragraph:

\textit{Radiation.}
Radiation transport is solved with the M1 moments method
\citep{Levermore_1984}, as implemented in {\small GIZMO} and
described in \citet{FORGE_d_in_FIRE_Hopkin_2023,STARFORGE_Towa_Grudic_2021},
across five frequency bands: H-ionizing, FUV/photo-electric, NUV,
optical--NIR, and an adaptive-wavelength FIR band that tracks the
evolving dust emission temperature.
Gas, dust, and radiation temperatures are evolved separately with
appropriate coupling terms, allowing self-consistent treatment of
both optically thin and optically thick regimes.
We use a reduced speed of light $\tilde{c} = 0.1\,c$
\note{[VERIFY: confirm \tilde{c} value used]},
following \citet{FORGE_d_in_FIRE_Hopkin_2023}.

%%  NB: if LW is a separate band in your sim (not collapsed into FUV),
%%  keep it listed. [VERIFY with simulation setup.]


%% ============================================================
%% 6.  THERMOCHEMISTRY: fill in the \note{TODO} on H2 dissociation
%% ============================================================
%%
%% CURRENT: "Thermodynamic contribution of H2 dissociation is not
%%   resolved in the code, however \note{TODO: why this is not a problem}"
%%
%% ISSUE: This note should be filled in. The standard argument is:
%%   (a) H2 dissociation heating (~4.5 eV/molecule) is a small correction
%%       to the total energy budget at densities where it occurs
%%       (n > 10^14 cm^-3, where compressional heating dominates).
%%   (b) The second collapse (at n ~ 10^16 cm^-3) driven by dissociation
%%       is not resolved since sinks form at n_thresh ~ 10^8.8 cm^-3.
%%   Meziani+2026 acknowledges this limitation and argues that since
%%   protostellar evolution tracks (Offner 2009) absorb this energy
%%   implicitly through the stellar radius evolution, it is not double-counted.
%%
%% SUGGESTED REPLACEMENT:

Thermodynamic energy release from H$_2$ dissociation
($\sim 4.5$~eV per molecule) is not explicitly tracked in the gas
energy equation.
This is acceptable because: (i) sink particles form at
$n_\mathrm{thresh} \approx 10^{8.8}\,\mathrm{cm}^{-3}$, well
before the second-collapse dissociation front
($n \gtrsim 10^{16}\,\mathrm{cm}^{-3}$) develops; and (ii)
the protostellar evolution tracks of \citet{Offner_2009}
self-consistently evolve the stellar radius (and hence the
protostellar luminosity and feedback) without requiring the gas
to track dissociation energy explicitly.


%% ============================================================
%% 7.  THERMOCHEMISTRY: add dust-to-gas ratio statement
%% ============================================================
%%
%% CURRENT: no mention of dust-to-gas ratio.
%% ISSUE: FORGE uses f_dg = 0.01*(Z/Zsun)*exp(...). For Pop III (Z=0),
%%   f_dg = 0 exactly, so dust plays no role. This should be stated to
%%   clarify that H2 formation on dust (the dominant channel in present-day
%%   ISM) is absent, and gas-phase channels dominate.
%%
%% ADD to Thermochemistry paragraph:

Since we simulate metal-free primordial gas, the dust-to-gas ratio
is zero and the dust term in the H$_2$ formation equation
(Eq.~\ref{eq:h2}) is absent.
H$_2$ formation proceeds exclusively via gas-phase channels:
the two-body H$^{-}$ and H$_2^+$ pathways at low densities
($n \lesssim 10^8\,\mathrm{cm}^{-3}$) and the three-body
reaction (3H $\to$ H$_2$ + H) at high densities.


%% ============================================================
%% 8.  THERMOCHEMISTRY: UV background self-shielding
%% ============================================================
%%
%% CURRENT: "UV Background is taken as tabulated values from
%%   \citep{Faucher-Giguère_2019} model."
%% ISSUE: Missing self-shielding. FORGE/FIRE-3 include self-shielding
%%   of H2 against LW photons and self-shielding of HI. At z~21
%%   the meta-galactic UV background is negligible, but local LW
%%   radiation from nearby protostars (resolved in radiation transport)
%%   is the dominant photodissociation source. This distinction matters.
%%
%% ADD:

At $z \sim 21$, the meta-galactic UV background
\citep{Faucher-Giguère_2019} is negligible compared to local
protostellar radiation resolved in the simulation.
H$_2$ self-shielding against Lyman--Werner photons is included
following the column-density approximation of
\note{[ADD: shielding function reference, e.g.\ Wolcott-Green et al. 2011]}.


%% ============================================================
%% 9.  COSMIC RAYS: add transport details
%% ============================================================
%%
%% CURRENT: "we evolve cosmic rays coupled with thermochemistry.
%%   Cosmic rays acceleration is one of potential stellar feedback
%%   mechanisms."
%% ISSUE: Very vague. FIRE-3 CR treatment includes streaming and
%%   diffusion transport. At z~21 with no SNe having occurred,
%%   CRs from pre-existing sources are negligible; only if stellar
%%   feedback launches jets/winds does CR injection become relevant.
%%   Should either describe the CR transport model or note CRs are
%%   sub-dominant in this regime.
%%
%% SUGGESTED REPLACEMENT:

\textit{Cosmic rays.}
Cosmic rays are evolved with streaming and diffusion transport
as described in \citet{FIRE_3_Updated_Hopkin}.
In our simulation, which covers only $\sim 15\kyr$ of post-collapse
evolution and does not include SNe, cosmic ray injection from
stellar feedback is negligible; the initial cosmic ray background
at $z \sim 21$ is at the level of the primordial cosmological value
and contributes subdominantly to H$_2$ dissociation and ionization
relative to the resolved radiation field.


%% ============================================================
%% 10. STELLAR EVOLUTION: fill in the \note{Confirm this!} for sink
%%     formation criterion
%% ============================================================
%%
%% CURRENT: "A gas particle is converted to a star particle once the
%%   density criterion is met \note{Confirm this! More info on the
%%   criterion}."
%%
%% ISSUE: The STARFORGE sink formation criterion (from FORGE paper
%%   Sec. 2.3 and Grudic+2021) is much more stringent than just density:
%%   it requires ALL of: (i) virial/self-gravity, (ii) Jeans unstable,
%%   (iii) converging flow, (iv) fully compressive tides in all directions,
%%   (v) local density/potential maximum, (vi) no pre-existing neighbouring
%%   sink within the accretion radius, (vii) collapse time < infall time
%%   onto nearest sink.
%%   [CRITICAL] Verify which of these criteria your simulation uses.
%%
%% SUGGESTED REPLACEMENT:

A gas particle is converted to a sink (protostar) particle when
it meets the full STARFORGE formation criteria
\citep{STARFORGE_Towa_Grudic_2021}: the cell must be
(i) locally self-gravitating (virial parameter $\alpha_\mathrm{vir} < 1$),
(ii) Jeans unstable,
(iii) in a converging flow ($\nabla \cdot \mathbf{v} < 0$),
(iv) experiencing compressive tides in all directions,
(v) a local maximum of the density field, and
(vi) not within the accretion radius of a pre-existing sink.
\note{[VERIFY: confirm all criteria are active; add any additional
criteria specific to your implementation.]}


%% ============================================================
%% 11. STELLAR EVOLUTION: fill in \note{TODO} section, add jets
%% ============================================================
%%
%% CURRENT: ends with "\note{TODO}"
%% ISSUE: STARFORGE protostars launch jets. Meziani+2026 is specifically
%%   about jet effects. Your sim runs 15 kyr -- do jets launch in this time?
%%   From the mu_phi discussion, jets should launch when mu_phi ~ few.
%%   [CRITICAL] State whether jets are active in this simulation or not,
%%   and at what point they might have activated.
%%
%% SUGGESTED:

Once formed, sink particles accrete bound gas whose circularised
orbit falls within the sink radius, following the accretion
model validated in \citet{STARFORGE_Towa_Grudic_2021}.
Each sink evolves along protostellar tracks following
\citet{Offner_2009}, explicitly tracking the stellar radius,
luminosity, and effective temperature as a function of the
instantaneous mass and accretion rate.
Protostellar jets are launched with a mass-loading proportional
to the surface accretion rate and a launch velocity comparable
to the escape velocity from the protostellar surface
\citep{STARFORGE_Towa_Grudic_2021}.
\note{[CRITICAL] Confirm whether jets are active in this simulation.
If jet feedback has not yet launched (expected when $\mu_\Phi \sim
\mathrm{few}$, which the simulation approaches by $\sim 10\kyr$),
state this explicitly.}


%% ============================================================
%% 12. INITIAL CONDITIONS: gravitational softening values
%% ============================================================
%%
%% CURRENT: "Stellar particles possess fixed softening lengths of
%%   \note{VALUE}."
%% ISSUE: FORGE uses Plummer-equivalent epsilon ~ 25 au (Delta x ~ 40 au)
%%   for sink particles, and adaptive softening for gas matching the
%%   hydrodynamic resolution. [CRITICAL] Fill in your actual value.
%%
%% SUGGESTED:

Gas particles use adaptive Lagrangian gravitational softening
matching the hydrodynamic kernel radius.
Sink/star particles use a fixed Plummer-equivalent softening
$\epsilon = $ \note{VALUE}~au, comparable to the accretion radius.
The fifth-order Hermite integrator of \citet{STARFORGE_Towa_Grudic_2021}
is used for all sink--sink and sink--gas gravitational encounters to
accurately resolve close binaries and dynamical ejections.


%% ============================================================
%% 13. REFINEMENT: substantially expand this subsection
%% ============================================================
%%
%% CURRENT: very sparse, multiple \note{} placeholders.
%% ISSUE: The refinement strategy is one of the most novel aspects of
%%   FORGE-style simulations and needs proper description.
%%   Key points from FORGE and Meziani+2026:
%%   - Jeans refinement: gas cells split when mass > f_J * M_Jeans,
%%     with f_J ~ 0.001 (i.e. ~1000 cells per Jeans mass)
%%   - Proximity refinement: additional refinement near sinks
%%   - Gradual ramp: each refinement level evolved for a few dynamical
%%     times before next level activates (avoids memory crashes and
%%     spurious initial conditions from geometry)
%%   [CRITICAL] Fill in your specific values of f_J, N_levels, r_refine.
%%
%% SUGGESTED REPLACEMENT for the Refinement subsection:

\subsection{Refinement strategy}
\label{sec:refinement}

We first evolve the full cosmological simulation with FIRE-3
physics until the gas in the target minihalo reaches a density
\note{$n \sim $ VALUE}~cm$^{-3}$, just before the first protostar
would form in the unrefined run.
We then restart from the preceding snapshot with an additional
``hyper-refinement'' layer applied in the neighbourhood of the
collapsing core.

Gas cells are refined (split into pairs of equal-mass daughter
cells) whenever their mass exceeds a fraction $f_J$ of the
local Jeans mass,
\begin{equation}
    m_\mathrm{cell} > f_J \, M_J \equiv f_J\,\frac{\pi^{5/2}}{6}
    \frac{c_s^3}{G^{3/2}\rho^{1/2}},
\end{equation}
with $f_J \approx 10^{-3}$ \note{[VERIFY]}, corresponding to
$\sim 1000$ resolution elements per Jeans mass and satisfying the
\citet{truelove97} criterion ($f_J \lesssim 10^{-2}$, or $\gtrsim
4$ cells per Jeans length) by a large margin.
Additional refinement is applied within a radius
$r_\mathrm{refine} = $ \note{VALUE}~pc of any sink particle,
progressively increasing resolution toward the sink.

To avoid spurious features from sudden cell-geometry imprinting
and memory overloads from creating large numbers of cells
simultaneously, the refinement is activated gradually: each
concentric ``shell'' in radius is evolved for a few local
dynamical times before the next inner shell is refined
\citep{FORGE_d_in_FIRE_Hopkin_2023}.
The full refinement sequence spans \note{N} levels from the
cosmological box down to the protostellar disk scale.

At our highest refinement level, the gas mass resolution is
$\Delta m \approx 10^{-3}\,\Msun$ and the spatial resolution
is $\Delta x \approx $ \note{VALUE}~AU, sufficient to resolve
densities up to $n \sim 10^{15}\,\mathrm{cm}^{-3}$ and
timescales down to $\sim $ \note{VALUE}~yr.


%% ============================================================
%% 14. CUTOUT: fill in the \note{Why?}
%% ============================================================
%%
%% CURRENT: "\note{Why? Idea: impact of LSS is not significant on the
%%   resolved timescales}"
%% SUGGESTED:

The cutout is used for analysis (not for re-simulation): over
the $\sim 15\kyr$ of post-collapse evolution studied here,
the sound-crossing time across the cutout radius
($r_\mathrm{cut} \approx 0.32$~pc) is
$\sim r_\mathrm{cut}/c_s \sim $ \note{VALUE}~kyr, comparable
to or longer than our integration time.
The large-scale structure outside the cutout therefore has no
causal influence on the disk dynamics on these timescales, and
the cutout fully captures the relevant physics.


%% ============================================================
%% 15. STRUCTURE NOTE: subsubsection vs subsection mismatch
%% ============================================================
%%
%% The "Hydrodynamical+Gravitational solver" is marked as \subsubsection
%% but all other solver descriptions are \subsection or below a \subsection
%% "Physics summary". This is structurally inconsistent -- the solver
%% description should be a \subsection (or the Physics summary should be
%% removed and its content merged into the solver/MHD/radiation subsections).
%%
%% SUGGESTED STRUCTURE:
%%
%% \section{Simulation and Methods}
%%   \subsection{Code and hydrodynamical solver}   (currently subsubsection)
%%   \subsection{Magnetic fields}
%%   \subsection{Radiation}
%%   \subsection{Thermochemistry}
%%   \subsection{Cosmic rays}
%%   \subsection{Stellar formation and evolution}
%%   \subsection{Initial conditions}
%%   \subsection{Refinement strategy}
%%   \subsection{Cutout re-simulation}
%%   \subsection{Disk identification}
%%   \subsection{Velocity dispersion and Toomre Q}


%% ============================================================
%% 16. MHD DIVERGENCE CONTROL: clarify cleaning vs constrained-gradient
%% ============================================================
%%
%% CURRENT: cites Hopkins_Raives_2015 for MHD. Item 2 above adds the
%%   constrained-gradient reference. But these are two different things:
%%
%%   Hopkins & Raives (2016, MNRAS 455, 51) = the MFM-MHD scheme itself,
%%     using hyperbolic/parabolic divergence cleaning (Powell + Dedner style).
%%     This is the DEFAULT in FIRE-2 and FIRE-3.
%%
%%   Hopkins (2016, MNRAS 462, 576) = the constrained-GRADIENT method,
%%     a separate improvement for greater divergence accuracy.
%%     This is what FORGE'd in FIRE explicitly uses.
%%
%% [CRITICAL] Which one does your simulation use? If you use GIZMO with
%%   the default FIRE-3 setup, you likely use the first (divergence cleaning).
%%   If you adopted the FORGE improvements, you use the second.
%%   Cite the correct paper. Do not conflate them.
%%
%%   If using divergence cleaning, the sentence should read:
%%
%%     The ideal MHD equations are integrated as described in
%%     \citet{Hopkins_Raives_2015} using a hyperbolic/parabolic
%%     divergence-cleaning scheme following \citet{Dedner_2002}.
%%
%%   If using constrained-gradient:
%%
%%     The ideal MHD equations are integrated as described in
%%     \citet{Hopkins_Raives_2015}, using the constrained-gradient
%%     method of \citet{Hopkins_2016_CG_MHD} for greater accuracy.

%% NEW BIB ENTRIES NEEDED:
%% @article{Hopkins_Raives_2015,   %% (check if already in bib as A_new_class_of_Hopkin_2015)
%%     author = {Hopkins, Philip F. and Raives, Matthew J.},
%%     title  = "{Accurate, meshless methods for magnetohydrodynamics}",
%%     journal= {\mnras},
%%     volume = {455},
%%     pages  = {51--88},
%%     year   = {2016},
%%     doi    = {10.1093/mnras/stv2180}
%% }
%%
%% @article{Dedner_2002,
%%     author = {Dedner, A. and Kemm, F. and Kr{\"o}ner, D. and Munz, C.-D.
%%               and Schnitzer, T. and Wesenberg, M.},
%%     title  = "{Hyperbolic Divergence Cleaning for the MHD Equations}",
%%     journal= {Journal of Computational Physics},
%%     volume = {175},
%%     pages  = {645--673},
%%     year   = {2002},
%%     doi    = {10.1006/jcph.2001.6961}
%% }


%% ============================================================
%% 17. THREE-BODY H2 FORMATION: cite explicitly
%% ============================================================
%%
%% CONFIRMED: Three-body H2 formation (3H → H2 + H) IS in FIRE-3.
%%   FIRE-3 Section 3.5 gives α_{3B} = 6×10^{-32} T^{-0.25} + 2×10^{-31} T^{-0.5}
%%   explicitly. This was previously incorrect in the suggestions list.
%%
%% The methods text mentions "alpha_{3B} term" but gives no citation.
%%   Standard references for the rate coefficient are:
%%   - Palla, Salpeter & Stahler (1983)
%%   - Abel, Bryan & Norman (2002)
%%
%% SUGGESTED addition:

The three-body formation coefficient $\alpha_{3\mathrm{B}}$
follows \citet{Abel_2002} and \citet{Palla_1983}; this reaction
becomes the dominant H$_2$ formation channel for
$n \gtrsim 10^8\,\mathrm{cm}^{-3}$ and is essential for
reproducing the collapse of primordial gas to protostellar densities.

%% NEW BIB ENTRIES NEEDED:
%% @article{Abel_2002,
%%     author  = {Abel, Tom and Bryan, Greg L. and Norman, Michael L.},
%%     title   = "{The Formation of the First Star in the Universe}",
%%     journal = {Science},
%%     volume  = {295},
%%     pages   = {93--98},
%%     year    = {2002},
%%     doi     = {10.1126/science.1063991}
%% }
%%
%% @article{Palla_1983,
%%     author  = {Palla, F. and Salpeter, E. E. and Stahler, S. W.},
%%     title   = "{Primordial star formation - The role of molecular hydrogen}",
%%     journal = {\apj},
%%     volume  = {271},
%%     pages   = {632--641},
%%     year    = {1983},
%%     doi     = {10.1086/161231}
%% }


%% ============================================================
%% 18. HD COOLING: cite it (it IS in FIRE-3)
%% ============================================================
%%
%% CONFIRMED: HD cooling IS in FIRE-3. FIRE-3 Section 3.5 explicitly
%%   includes Λ_{HD} from Galli & Palla (1998) in the Λ_neutral term.
%%   This was previously listed as missing — that was incorrect.
%%   Your methods section already lists Λ_{HD}. Just add the citation.
%%   HD cooling dominates below T ~ 200 K (the H2 loitering point)
%%   and can allow lower fragmentation masses in metal-free gas.
%%
%% Standard references for HD cooling:
%%   - Flower et al. (2000, MNRAS 314, 753) -- HD-H collision rates
%%   - Nakamura & Umemura (2002, ApJ 569, 549) -- role in Pop III
%%   - Bromm & Larson (2004, ARA&A 42, 79) -- review
%%
%% SUGGESTED addition to Thermochemistry:

HD cooling ($\Lambda_\mathrm{HD}$), which dominates below
$T \lesssim 200$~K \note{[VERIFY whether HD is active at these
densities in your sim]}, is included following
\citet{Flower_2000} \note{[add correct rate reference]}.


%% ============================================================
%% 19. H- GAS-PHASE H2 PATHWAY: cite the reaction network
%% ============================================================
%%
%% Your equation includes \alpha_{GP} (gas-phase formation via H-/H2+
%%   pathways). The reference given is Glover_Jappsen_2007. This is
%%   correct but incomplete -- the H- pathway was tabulated by:
%%   - Abel et al. (1997, NewA 2, 181) -- H- rate coefficients
%%   - Galli & Palla (1998, A&A 335, 403) -- comprehensive primordial network
%%
%% SUGGESTED: supplement the Glover & Jappsen citation:

Gas-phase H$_2$ formation via the H$^{-}$ and H$_2^+$ pathways
($\alpha_\mathrm{GP}$) follows \citet{Glover_Jappsen_2007},
with rate coefficients from \citet{Abel_1997} and
\citet{Galli_Palla_1998}.

%% NEW BIB ENTRIES:
%% @article{Galli_Palla_1998,
%%     author  = {Galli, D. and Palla, F.},
%%     title   = "{The chemistry of the early Universe}",
%%     journal = {A\&A},
%%     volume  = {335},
%%     pages   = {403--420},
%%     year    = {1998}
%% }
%%
%% @article{Abel_1997,
%%     author  = {Abel, T. and Anninos, P. and Zhang, Y. and Norman, M. L.},
%%     title   = "{Modeling primordial gas in numerical cosmology}",
%%     journal = {New Astronomy},
%%     volume  = {2},
%%     pages   = {181--207},
%%     year    = {1997},
%%     doi     = {10.1016/S1384-1076(97)00010-9}
%% }


%% ============================================================
%% 20. UV BACKGROUND: confirm FG2019 vs FG2020
%% ============================================================
%%
%% FIRE-2 uses Faucher-Giguère et al. (2009).
%% FIRE-3 updates to Faucher-Giguère (2020) -- published online ~2019,
%%   which is why it is sometimes cited as "2019" or "2020".
%% Your bib cites \citet{Faucher-Giguère_2019}.
%%
%% [VERIFY] Make sure this is the 2020 paper (MNRAS 493, 1614),
%%   not the older FG09 background. The correct citation should be:
%%   Faucher-Giguère (2020, MNRAS 493, 1614), arXiv:1901.09283.
%%   At z~21 the UV background is negligible regardless, but the
%%   correct citation matters for reproducibility.


%% ============================================================
%% 21. COLLISIONAL H2 DISSOCIATION: cite rate coefficients
%% ============================================================
%%
%% Your H2 equation has the last term summing over collisional
%%   dissociation by e-, H+, H2, HI, HeI, HeII. These rate coefficients
%%   need citations.
%%   Standard source: Glover & Abel (2008, MNRAS 388, 1627) or
%%   Martin, Schwarz & Mandy (1996) for H + H2 dissociation.
%%
%% SUGGESTED addition after the sum term:

Collisional dissociation rate coefficients $\beta_{\mathrm{H_2},i}$
for each species $i$ are taken from \citet{Glover_Abel_2008}.

%% NEW BIB ENTRY:
%% @article{Glover_Abel_2008,
%%     author  = {Glover, Simon C. O. and Abel, Tom},
%%     title   = "{Uncertainties in H2 and HD cooling functions for
%%                 primordial gas}",
%%     journal = {\mnras},
%%     volume  = {388},
%%     pages   = {1627--1651},
%%     year    = {2008},
%%     doi     = {10.1111/j.1365-2966.2008.13224.x}
%% }


%% ============================================================
%% 22. GRAVITATIONAL SOFTENING: cite GIZMO gravity paper
%% ============================================================
%%
%% Currently cites only FORGE for the 5th-order Hermite integrator.
%%   The underlying GIZMO gravity solver (TreePM) should also be cited:
%%   Hopkins (2015, MNRAS 450, 53) -- the main GIZMO paper.
%%   This is distinct from the MHD paper.
%%
%% SUGGESTED:

Gravity is solved with the TreePM algorithm of {\small GIZMO}
\citep{A_new_class_of_Hopkin_2015}, using adaptive Lagrangian
force softening for gas and fixed softening for stellar/sink
particles, with a fifth-order Hermite integrator for close
encounters \citep{FORGE_d_in_FIRE_Hopkin_2023}.


%% ============================================================
%% 24. FULL LIST OF MISSING .bib ENTRIES (summary)
%% ============================================================
%%
%% Hopkins 2016 constrained-gradient MHD:
%% @article{Hopkins_2016_CG_MHD,
%%     author = {Hopkins, Philip F.},
%%     title  = "{A constrained-gradient method to control divergence errors in numerical MHD}",
%%     journal= {\mnras},
%%     volume = {462},
%%     pages  = {576--587},
%%     year   = {2016},
%%     doi    = {10.1093/mnras/stw1578}
%% }
%%
%% Wolcott-Green et al. 2011 (H2 self-shielding):
%% @article{wolcott_green11,
%%     author = {Wolcott-Green, Jemma and Haiman, Zolt\'{a}n and Bryan, Greg L.},
%%     title  = "{Photodissociation of H$_2$ in protogalaxies: modelling self-shielding in three-dimensional simulations}",
%%     journal= {\mnras},
%%     volume = {418},
%%     pages  = {838--852},
%%     year   = {2011},
%%     doi    = {10.1111/j.1365-2966.2011.19538.x}
%% }
%%
%% Glover & Jappsen 2007 (gas-phase H2 formation, already cited as
%%   Glover_Jappsen_2007 -- keep)
%%
%% Offner et al. 2009 (protostellar tracks, already cited as Offner_2009):
%% @article{Offner_2009,
%%     author  = {Offner, Stella S. R. and Klein, Richard I. and McKee, Christopher F. and Krumholz, Mark R.},
%%     title   = "{The Effects of Radiative Transfer on Low-Mass Star Formation}",
%%     journal = {\apj},
%%     volume  = {703},
%%     pages   = {131--149},
%%     year    = {2009},
%%     doi     = {10.1088/0004-637X/703/1/131}
%% }

%% ============================================================
%% END OF SUGGESTIONS
%% ============================================================
